% GENERATED FILE. Edit canonical lessons, then run npm run generate:books.
% canonical-source-hash: fnv1a64:f53f5e772a8ee473
% canonical-lessons: GU-C159-majvu, GU-C159-sukavvu, GU-C159-nal, GU-C159-dholvu, GU-C159-bhoy

\chapter{Scouring, drying and spills}
\label{ch:gu-scouring-drying-and-spills}
\hlchaptermodality{159}

\begin{chapteropening}
\textbf{By the end of this chapter:} \emph{I can talk about washing up, drying clothes and spills.}

Complete the last lesson of chapter 159: I can talk about washing up, drying clothes and spills.
\end{chapteropening}

\section[\texorpdfstring{\gu{માંજવું} (mā̃jvũ)}{માંજવું (mā̃jvũ)}]{\gu{માંજવું} (mā̃jvũ) — to scour}

\label{lesson:GU-C159-majvu}



\begin{quote}
Before the new one: say the Gujarati for floor cloth, then the Gujarati for stain.
\end{quote}

\subsection*{You'll want to know: \gu{માંજવું}}
\textbf{\gu{માંજવું}} — \emph{mā̃jvũ} — \textquotedblleft{}to scour\textquotedblright{}.

\begin{quote}
\textbf{\gu{તે} \gu{રાત્રે} \gu{વાસણ} \gu{માંજે} \gu{છે}.} — \emph{te rātre vāsaṇ mā̃je che} — \textquotedblleft{}She scours the pots at night.\textquotedblright{}
\end{quote}

\begin{cousinweb}
\textbf{\gu{માંજવું}} (\emph{mā̃jvũ}) is a tadbhava. It comes from Sanskrit \emph{mārjati}, \textquotedblleft{}wipes off\textquotedblright{}, on the root that gives \emph{mṛjanti}, \textquotedblleft{}they rub, clean\textquotedblright{}; CDIAL glosses the Gujarati verb \textquotedblleft{}to scour\textquotedblright{} (CDIAL 10080).
\end{cousinweb}

\begin{grammarlens}[title={What you've built}]
One more word for everyday life at A2.
\end{grammarlens}

\subsection*{Your turn}
\begin{itemize}
\raggedright
  \item \emph{Say it:} \emph{mā̃jvũ}
  \item \emph{Say it:} \emph{mā̃jvũ}, once more
  \item \emph{Say it:} \emph{te rātre vāsaṇ mā̃je che} — She scours the pots at night.
  \item \emph{Recall:} say the Gujarati for floor cloth, then the Gujarati for stain, then say \emph{mā̃jvũ} again
\end{itemize}

\begin{tcolorbox}[breakable,colback=teal!4,colframe=teal!35!black,title={Before you move on}]
What does \gu{માંજવું} mean? (\textquotedblleft{}To scour\textquotedblright{}.) Say it once more.
\end{tcolorbox}
\section[\texorpdfstring{\gu{સૂકવવું} (sūkavvũ)}{સૂકવવું (sūkavvũ)}]{\gu{સૂકવવું} (sūkavvũ) — to dry out, to put out to dry}

\label{lesson:GU-C159-sukavvu}



\begin{quote}
Before the new one: say the Gujarati for stain, then the Gujarati for to scour.
\end{quote}

\subsection*{You'll want to know: \gu{સૂકવવું}}
\textbf{\gu{સૂકવવું}} — \emph{sūkavvũ} — \textquotedblleft{}to dry out, to put out to dry\textquotedblright{}.

\begin{quote}
\textbf{\gu{માતા} \gu{અગાસીમાં} \gu{કપડાં} \gu{સૂકવે} \gu{છે}.} — \emph{mātā agāsīmā̃ kapḍā̃ sūkave che} — \textquotedblleft{}Mother dries the clothes on the terrace.\textquotedblright{}
\end{quote}

Related to \textbf{\gu{સૂકું}}, \emph{sūkũ}, dry.

\begin{cousinweb}
\textbf{\gu{સૂકવવું}} (\emph{sūkavvũ}) is a tadbhava, the causative of an inherited verb for drying. CDIAL lists Gujarati \emph{sukvũ}, \textquotedblleft{}to dry up\textquotedblright{}, and its causative \emph{sukavvũ} under \emph{*śuṣkati}, \textquotedblleft{}becomes dry\textquotedblright{} (CDIAL 12552).
\end{cousinweb}

\begin{grammarlens}[title={What you've built}]
One more word for everyday life at A2.
\end{grammarlens}

\subsection*{Your turn}
\begin{itemize}
\raggedright
  \item \emph{Say it:} \emph{sūkavvũ}
  \item \emph{Say it:} \emph{sūkavvũ}, once more
  \item \emph{Say it:} \emph{mātā agāsīmā̃ kapḍā̃ sūkave che} — Mother dries the clothes on the terrace.
  \item \emph{Recall:} say the Gujarati for stain, then the Gujarati for to scour, then say \emph{sūkavvũ} again
\end{itemize}

\begin{tcolorbox}[breakable,colback=teal!4,colframe=teal!35!black,title={Before you move on}]
What does \gu{સૂકવવું} mean? (\textquotedblleft{}To dry out, to put out to dry\textquotedblright{}.) Say it once more.
\end{tcolorbox}
\section[\texorpdfstring{\gu{નળ} (naḷ)}{નળ (naḷ)}]{\gu{નળ} (naḷ) — water tap, pipe}

\label{lesson:GU-C159-nal}



\begin{quote}
Before the new one: say the Gujarati for to scour, then the Gujarati for to dry out.
\end{quote}

\subsection*{You'll want to know: \gu{નળ}}
\textbf{\gu{નળ}} — \emph{naḷ} — \textquotedblleft{}water tap, pipe\textquotedblright{}.

\begin{quote}
\textbf{\gu{રસોડામાં} \gu{નળ} \gu{છે}.} — \emph{rasoḍāmā̃ naḷ che} — \textquotedblleft{}There is a tap in the kitchen.\textquotedblright{}
\end{quote}

\begin{cousinweb}
\textbf{\gu{નળ}} (\emph{naḷ}) is a tadbhava. It comes from Sanskrit \emph{naḍa}, \emph{naḷa}, \textquotedblleft{}a kind of reed\textquotedblright{}, a hollow stem; CDIAL glosses the Gujarati word \textquotedblleft{}pipe, drain\textquotedblright{} (CDIAL 6936).
\end{cousinweb}

\begin{grammarlens}[title={What you've built}]
One more word for everyday life at A2.
\end{grammarlens}

\subsection*{Your turn}
\begin{itemize}
\raggedright
  \item \emph{Say it:} \emph{naḷ}
  \item \emph{Say it:} \emph{naḷ}, once more
  \item \emph{Say it:} \emph{rasoḍāmā̃ naḷ che} — There is a tap in the kitchen.
  \item \emph{Recall:} say the Gujarati for to scour, then the Gujarati for to dry out, then say \emph{naḷ} again
\end{itemize}

\begin{tcolorbox}[breakable,colback=teal!4,colframe=teal!35!black,title={Before you move on}]
What does \gu{નળ} mean? (\textquotedblleft{}Water tap, pipe\textquotedblright{}.) Say it once more.
\end{tcolorbox}
\section[\texorpdfstring{\gu{ઢોળવું} (ḍhoḷvũ)}{ઢોળવું (ḍhoḷvũ)}]{\gu{ઢોળવું} (ḍhoḷvũ) — to spill, to pour out}

\label{lesson:GU-C159-dholvu}



\begin{quote}
Before the new one: say the Gujarati for to dry out, then the Gujarati for water tap.
\end{quote}

\subsection*{You'll want to know: \gu{ઢોળવું}}
\textbf{\gu{ઢોળવું}} — \emph{ḍhoḷvũ} — \textquotedblleft{}to spill, to pour out\textquotedblright{}.

\begin{quote}
\textbf{\gu{બાળક} \gu{દૂધ} \gu{ઢોળે} \gu{છે}.} — \emph{bāḷak dūdh ḍhoḷe che} — \textquotedblleft{}The child spills the milk.\textquotedblright{}
\end{quote}

\begin{cousinweb}
\textbf{\gu{ઢોળવું}} (\emph{ḍhoḷvũ}) is a tadbhava. It comes from \emph{*ḍhōlayati}, the causative of \emph{*ḍhulati}, \textquotedblleft{}bends, falls, flows\textquotedblright{}, a reconstructed form; CDIAL glosses the Gujarati verb \textquotedblleft{}to pour out\textquotedblright{} (CDIAL 5593).
\end{cousinweb}

\begin{grammarlens}[title={What you've built}]
One more word for everyday life at A2.
\end{grammarlens}

\subsection*{Your turn}
\begin{itemize}
\raggedright
  \item \emph{Say it:} \emph{ḍhoḷvũ}
  \item \emph{Say it:} \emph{ḍhoḷvũ}, once more
  \item \emph{Say it:} \emph{bāḷak dūdh ḍhoḷe che} — The child spills the milk.
  \item \emph{Recall:} say the Gujarati for to dry out, then the Gujarati for water tap, then say \emph{ḍhoḷvũ} again
\end{itemize}

\begin{tcolorbox}[breakable,colback=teal!4,colframe=teal!35!black,title={Before you move on}]
What does \gu{ઢોળવું} mean? (\textquotedblleft{}To spill, to pour out\textquotedblright{}.) Say it once more.
\end{tcolorbox}
\section[\texorpdfstring{\gu{ભોંય} (bhõy)}{ભોંય (bhõy)}]{\gu{ભોંય} (bhõy) — floor, ground}

\label{lesson:GU-C159-bhoy}



\begin{quote}
Before the new one: say the Gujarati for water tap, then the Gujarati for to spill.
\end{quote}

\subsection*{You'll want to know: \gu{ભોંય}}
\textbf{\gu{ભોંય}} — \emph{bhõy} — \textquotedblleft{}floor, ground\textquotedblright{}.

\begin{quote}
\textbf{\gu{ભોંય} \gu{ભીની} \gu{છે}.} — \emph{bhõy bhīnī che} — \textquotedblleft{}The floor is wet.\textquotedblright{}
\end{quote}

\begin{cousinweb}
\textbf{\gu{ભોંય}} (\emph{bhõy}) is a tadbhava. It comes from Sanskrit \emph{bhūmi}, \textquotedblleft{}earth, ground\textquotedblright{}, and CDIAL lists Gujarati \emph{bhõ} and \emph{bhõy}, \textquotedblleft{}ground\textquotedblright{}, marking \emph{bhõy} as the Kathiawar form (CDIAL 9557).
\end{cousinweb}

\begin{grammarlens}[title={What you've built}]
One more word for everyday life at A2.
\end{grammarlens}

\subsection*{Your turn}
\begin{itemize}
\raggedright
  \item \emph{Say it:} \emph{bhõy}
  \item \emph{Say it:} \emph{bhõy}, once more
  \item \emph{Say it:} \emph{bhõy bhīnī che} — The floor is wet.
  \item \emph{Recall:} say the Gujarati for water tap, then the Gujarati for to spill, then say \emph{bhõy} again
\end{itemize}

\begin{tcolorbox}[breakable,colback=teal!4,colframe=teal!35!black,title={Before you move on}]
What does \gu{ભોંય} mean? (\textquotedblleft{}Floor, ground\textquotedblright{}.) Say it once more.
\end{tcolorbox}
